\section{Method}
\label{sec:method}

Figure~\ref{fig:architecture} gives an overview of MedAlign-RAG. A query is
served by two retrieval paths whose budget is allocated by a density-aware
router; the mixed candidate pool is ordered by an off-the-shelf reranker and
then passed to a vision-language agent that judges the value of each candidate
and either accepts the evidence for generation or rewrites the query and
re-retrieves. The agent is trained by reinforcement learning against an
answer-utility reward. Once accepted, the selected image--text evidence is fed,
unaltered, to a vision-language generator.

\begin{figure*}[t]
  \centering
  \IfFileExists{figures/fig2_architecture.pdf}{%
    \includegraphics[width=0.95\textwidth]{fig2_architecture.pdf}}{%
    \framebox[0.95\textwidth]{\rule{0pt}{4.2cm}\ Figure~2 placeholder (system architecture)}}
  \caption{MedAlign-RAG. Dual-path retrieval feeds a density-aware router that
  allocates the text/image budget; the reranked pool is inspected by a
  vision-language agent that judges each image and text candidate and either
  accepts the evidence for generation or rewrites the query and re-retrieves
  (closing the loop). The agent is trained by reinforcement learning against an
  answer-utility reward---the lift in the frozen generator's probability of the
  correct answer. Accepted image--text evidence is passed unaltered to the
  generator.}
  \label{fig:architecture}
\end{figure*}

\subsection{Problem Formulation and Evaluation Protocol}
\label{sec:formulation}

A query is a pair $x=(q, I_q)$, where $q$ is the question text (with answer
options for multiple-choice QA) and $I_q$ is an optional query image. Given a
text corpus $\mathcal{T}$ and an image corpus $\mathcal{V}$, a retriever returns
a candidate pool $\mathcal{C}=\mathcal{C}_t\cup\mathcal{C}_i$ with
$\mathcal{C}_t\subseteq\mathcal{T}$ and $\mathcal{C}_i\subseteq\mathcal{V}$. An
agent selects an evidence set $E\subseteq\mathcal{C}$---not in one shot, but
through iterative rounds of judging, generating, or reformulating and
re-retrieving (Sec.~\ref{sec:selection})---and a vision-language generator $G$
produces the answer
\begin{equation}
  \hat{a} = G\!\left(q, I_q, E\right).
  \label{eq:generation}
\end{equation}

The distinction we stress is the \emph{form} of $E$ passed to $G$. Prior work
adopts a \emph{text-proxy} protocol in which every image evidence $e\in E$ is
replaced by its textual caption $\text{cap}(e)$, so $G$ receives text only.
We instead adopt an \emph{end-to-end image--text} protocol in which the actual
pixels of image evidence are passed to $G$. As we show empirically, the two
protocols rank methods differently: text-proxy evaluation is blind to the
quality of visual evidence and therefore misestimates reranking quality on
image-dependent queries.

\subsection{Dual-Path Retrieval}
\label{sec:retrieval}

The text path encodes $q$ with a text embedding model~\cite{chen2024bgem3} and
searches $\mathcal{T}$ by inner product, returning the top $N_t$ passages. The
image path encodes $I_q$ with the vision tower of a vision-language
model~\cite{bai2023qwenvl} and searches $\mathcal{V}$, returning the top $N_i$
images. Each candidate $c$ carries a first-stage score $s(c)$ from its own path.
The union $\mathcal{C}$, with $|\mathcal{C}|\le N_t+N_i$, forms the input to
reranking.

\subsection{Density-Aware Modality Routing}
\label{sec:routing}

The retrieval budget should follow where a query's discriminative information
lives---and, just as importantly, avoid a modality that only injects noise. A
query phrased in generic or positional terms (\emph{``which anatomical site is
shown''}) carries no discriminative textual signal; issued to the text path it
retrieves same-site, different-pathology passages that flood the pool with
homogeneous false positives. Routing must therefore be able not only to
rebalance the two paths but to \emph{suppress} one of them.

We predict two properties of $q$ with two lightweight classifiers sharing a
Chinese BERT backbone. A \emph{textual density} head outputs
$d=f_{\text{den}}(q)\in\{L_0,\dots,L_4\}$, from $L_0$ (vague, few localizable
cues) to $L_4$ (highly specific, multiple explicit constraints). An
\emph{image-dependency} head outputs $r=f_{\text{dep}}(q)\in\{R_1,R_2,R_3\}$,
from $R_1$ (text usually suffices) to $R_3$ (the image is indispensable).

Given a total budget $B=N_t+N_i$, the router allocates the image share through a
monotone function $\phi$ that increases with image dependency and decreases with
textual density,
\begin{equation}
  N_i = \big\lceil B\cdot\phi(d, r)\big\rceil,\qquad N_t = B - N_i,
  \label{eq:budget}
\end{equation}
and produces a fusion weight $w\in[0,1]$ that biases the reranker toward the
trusted modality. Routing is re-run on every retrieval round, so a query
reformulated by the agent (Sec.~\ref{sec:selection}) is re-budgeted from its new
text. Crucially, $\phi$ saturates: a
low-density, image-critical query ($L_0,R_3$) drives $N_t\!\to\!0$ and
$w\!\to\!1$, gating off the noisy text path, whereas a text-answerable query
($L_4,R_1$) drives $N_i\!\to\!0$, avoiding spurious visual matches. Unlike a
fixed fusion weight, routing can thus withdraw a modality entirely when its
signal is noise rather than evidence.

\subsection{Agentic Cross-Modal Evidence Selection}
\label{sec:selection}

Rather than take the reranker's top $K$ as final, we treat evidence selection as
a sequential decision problem solved by a vision-language agent $\pi$. The agent
sees the actual pixels of image candidates and the text of passage candidates,
judges which ones are worth keeping, and---when the pool is inadequate---rewrites
the query and retrieves again. It is trained so that the evidence it keeps
maximizes the generator's chance of answering correctly, rather than surface
relevance.

\paragraph{Initial ordering.}
The mixed pool $\mathcal{C}$ is first ordered by an off-the-shelf multimodal
reranker, using the fusion weight $w$ from routing to bias the ranking toward the
trusted modality. This ordering is only a prior; the agent, not the reranker,
decides what is used.

\paragraph{Agent policy.}
The policy $\pi_\theta$ is a small vision-language model that ingests the query
$(q,I_q)$ and the ranked candidates---image candidates as pixels, so visual
content is judged directly rather than through captions---and, at each round
$t$, emits an action $a_t$:
\begin{equation}
  a_t \sim \pi_\theta\!\left(\cdot \mid q, I_q, \mathcal{C}^{(t)}, E^{(t)}\right),
  \label{eq:policy}
\end{equation}
where $a_t$ either (i)~\emph{selects} candidates into the running evidence set
$E^{(t)}$, (ii)~\emph{accepts} $E^{(t)}$ and stops, or (iii)~\emph{rewrites} the
query to $q^{(t+1)}$, triggering a fresh retrieval round (through routing,
Sec.~\ref{sec:routing}) that refreshes the pool $\mathcal{C}^{(t+1)}$. Rewriting
gives the system a recourse absent from one-shot rerankers: a pool of
same-site, different-pathology false positives can be discarded and replaced by
a more discriminative query. The loop runs until the agent accepts or a round
budget $T$ is reached.

\paragraph{Answer-utility reward.}
The agent is rewarded by how much the accepted evidence improves the frozen
generator's confidence in the correct answer. For a multiple-choice query with
gold answer token $a^\star$, let $P_G(a^\star \mid q, I_q, E)$ be the generator's
probability of $a^\star$ given evidence $E$. The reward is the utility
\emph{lift} over answering with no evidence,
\begin{equation}
  R(E) = P_G\!\left(a^\star \mid q, I_q, E\right)
       - P_G\!\left(a^\star \mid q, I_q, \varnothing\right).
  \label{eq:reward}
\end{equation}
The second term is a per-query baseline that measures how much the generator
already knows closed-book, so the agent is credited only for evidence that
genuinely helps. Because $R$ is read from the generator's answer distribution, it
needs \emph{no passage-level relevance labels}---the multiple-choice answer is
the only supervision. Evidence that is topically related but useless yields no
lift, and confidently-wrong homogeneous-image matches, which pull probability
toward a distractor, are penalized---directly targeting~(P2).

\paragraph{Training.}
We optimize $\theta$ with policy-gradient reinforcement learning, maximizing the
expected reward $\mathbb{E}_{a_{1:T}\sim\pi_\theta}\!\left[R(E)\right]$ over
sampled trajectories, with the generator $G$ frozen throughout so the reward
reflects evidence quality alone. Because the reward is defined on the end-to-end
image--text answer distribution, training \emph{requires} the corrected
evaluation protocol (Sec.~\ref{sec:formulation}): under a text-proxy protocol the
generator never sees image evidence, and the utility signal that drives learning
would be blind to it.

\subsection{Multimodal Generation}
\label{sec:generation}

The selected set $E$---original image pixels together with text passages---is
provided to the vision-language generator following Eq.~\ref{eq:generation}. No
image is replaced by a caption, so visual evidence contributes directly to the
answer.
